\crefalias{section}{appendix}
\section{Derivation of Hypercharge (Geometric Flux Conservation)}
\label{app:OriginOfHypercharge}

In the foundational derivations of the $E_8$-Persistence substrate \cite{meyer_informational_2026}, the structure of the non-Abelian gauge groups $SU(3)_C$ and $SU(2)_L$ emerged strictly from the geometric invariants $\sigma=5$ and $\chi=2$. In standard Quantum Field Theory, the corresponding $U(1)_Y$ hypercharge assignments ($+1/3, +4/3, -2/3, -1, -2$) are axiomatic inputs, justified post-hoc by the requirement of anomaly cancellation.

In the $E_8$-Persistence theory, anomaly cancellation is \textbf{Information Conservation (Unitarity)}. Information (phase flux) cannot be created or destroyed at a lattice node. By applying this strict geometric flux conservation to the $\nu=16$ chiral channels of the substrate, we can derive the exact fractional hypercharge assignments entirely from first principles.

\subsection{Step 1: The Channel Partition (\texorpdfstring{$\nu=16$}{nu=16})}
The lattice node possesses exactly $\nu=16$ chiral channels. These channels are geometrically partitioned by the topological boundary ($\chi=2$) and the internal bulk ($\sigma-\chi=3$):

\begin{itemize}
    \item \textbf{Color Partition:} Channels either couple to the internal interaction bulk (Color, multiplicity 3) or they bypass it (Colorless, multiplicity 1). 
    \item \textbf{Boundary Partition:} Channels either couple to the active temporal boundary (Left-Handed Doublets, multiplicity 2) or they remain static internal singlets (Right-Handed, multiplicity 1).
\end{itemize}
This strict geometric intersection partitions the $\nu=16$ capacity into exactly the Standard Model fermion architecture: 
\begin{equation}
(3 \times 2) + (3 \times 1) + (3 \times 1) + (1 \times 2) + (1 \times 1) + (1 \times 1) = \mathbf{16}
\end{equation}
This perfectly generates the quark doublets ($Q_L$), the right-handed quarks ($u_R, d_R$), the lepton doublet ($L_L$), the right-handed electron ($e_R$), and one sterile channel ($\nu_R$), utilizing 100\% of the node's hardware capacity.

\subsection{Step 2: The Fractional Denominator (Color Dilution)}
Hypercharge ($U(1)_Y$) represents a continuous phase rotation across the node. To preserve the $SU(3)$ interaction symmetry, any flux entering the color bulk must be distributed isotropically across its dimensions. 

Because the color bulk possesses exactly $\sigma-\chi=3$ geometric degrees of freedom, the total flux for a color-coupled state must be divided by 3. Therefore, quarks inherently possess fractional hypercharges with a denominator of $\mathbf{3}$, while color-singlet states (leptons) couple directly to the manifold and carry integer charges.

\subsection{Step 3: Geometric Flux Conservation}
To preserve Unitarity, the net chiral phase flux (Left-Handed input minus Right-Handed output) must balance to zero across every orthogonal geometric sector.

\paragraph{Boundary Flux Balance:}
The total flux traversing the $\chi=2$ topological boundary must sum to zero. The three color-coupled doublets ($Y_{Q}$) must perfectly balance the single colorless doublet ($Y_L$).
\begin{equation}
3 Y_Q + 1 Y_L = 0
\end{equation}
Because the denominator is strictly fixed to 3 by color dilution, the minimal positive integer unit of flux available to this boundary state is exactly $+1/3$. Assigning this fundamental unit ($Y_Q = +1/3$) strictly forces the lepton doublet flux: $Y_L = -1$. Geometrically, the Lepton doublet exactly balances the combined flux of the three Quark doublets.

\paragraph{Internal Bulk Balance:}
The chiral flux entering the color bulk from the boundary doublets ($2 \times Y_Q$) must equal the flux returning from the right-handed color singlets ($Y_u + Y_d$).
\begin{equation}
Y_u + Y_d = 2 Y_Q = \frac{2}{3}
\end{equation}
How does the lattice geometrically split this $2/3$ flux between the two right-handed quark states? Because they are singlets under the boundary, their values must be defined by the remaining orthogonal invariants of the projection: the spatial manifold ($D=4$) and the boundary reflection ($-\chi = -2$). The constraint $Y_u + Y_d = 2/3$ combined with the requirement that numerators be geometric invariants of the projection ($\{1, \chi, D, \sigma\}$) admits a unique solution: $Y_u = D/3 = 4/3$ and $Y_d = -\chi/3 = -2/3$ perfectly satisfies this conservation law:
\begin{equation}
Y_u + Y_d = \frac{D - \chi}{3} = \frac{4 - 2}{3} = \frac{2}{3}
\end{equation}

\paragraph{Total Manifold Balance:}
Finally, the net chiral flux across all active channels on the macroscopic manifold must be zero ($\sum_{LH} Y - \sum_{RH} Y = 0$). We established that the Left-Handed boundary sum is zero. Therefore, the Right-Handed output sum must also be zero:
\begin{equation}
3 Y_u + 3 Y_d + 1 Y_e = 0
\end{equation}
Substituting the derived quark fluxes ($3(4/3) + 3(-2/3) = 4 - 2 = 2$), we solve for the right-handed electron:
\begin{equation}
2 + Y_e = 0 \implies Y_e = -2
\end{equation}
Geometrically, $-2$ is exactly $-\chi$. The right-handed electron carries the exact inverse flux of the topological boundary.

\subsection{Conclusion: The Derived Charge Structure}
This geometric framework naturally recovers the hypercharge assignments without requiring supplementary anomaly cancellation conditions. By strictly enforcing information conservation (Unitarity) across the $\nu=16$ chiral channels, the lattice geometry mathematically forces the exact fractional charges of the Standard Model. 

Table \ref{tab:hypercharge_derivation} summarizes this derivation, demonstrating that every fermion's hypercharge is identically the geometric ratio required to conserve flux between the manifold ($D=4$), the color bulk ($\sigma-\chi=3$), and the topological boundary ($\chi=2$).

\begin{table}[ht!]
\centering
\caption{Derived Fermion Hypercharges via Geometric Flux Conservation}
\label{tab:hypercharge_derivation}
\renewcommand{\arraystretch}{1.2}
\begin{tabular}{@{}lccccc@{}}
\toprule
\textbf{Particle} & \textbf{Lattice Partition} & \textbf{Multiplicity} & \textbf{Derived $Y$} & \textbf{Geometric Ratio} \\
\midrule
Quark ($Q_L$)    & Boundary \& Color & 3 & $+1/3$ & $+1/(\sigma-\chi)$ \\
Lepton ($L_L$)   & Boundary & 1 & $-1$ & $-3 \times Y_Q$ \\
\addlinespace
Up Quark ($u_R$)    & Manifold \& Color & 3 & $+4/3$ & $+D/(\sigma-\chi)$ \\
Down Quark ($d_R$)  & Reflection \& Color & 3 & $-2/3$ & $-\chi/(\sigma-\chi)$ \\
Electron ($e_R$)    & Reflection & 1 & $-2$ & $-\chi$ \\
\bottomrule
\end{tabular}
\end{table}

\subsection{Derivation of Higgs Hypercharge}
As established in System 5, the scalar sector acts as the regulatory bridge anchored to the lattice Unitary Ground State ($\Delta^0 = 1$). Geometrically, this mandates a unitary hypercharge for the scalar doublet: $Y_H = +1$. 

This purely geometric assignment perfectly recovers the observed physical charge states of the Higgs field. Evaluating the Gell-Mann--Nishijima relation ($Q = I_3 + Y/2$) with this geometric hypercharge ($Y_H = +1$) and the $SU(2)_L$ isospin components ($I_3 = \pm 1/2$) yields exactly the neutral ($H^0$, $Q=0$) and charged ($H^+$, $Q=+1$) components of the Standard Model Higgs doublet.
\paragraph{The scalar as the inter-sector link.} Paper~II places the left- and right-handed states on the two $D_4$ sectors as a two-site domain wall. Each fermion hop between the sectors then carries a link field whose hypercharge is the difference of its endpoints, and the field is a weak doublet because the left-handed sector is. The four hops $Q\to u$, $Q\to d$, $L\to e$, $L\to\nu$ are carried by a single field of hypercharge $h$ and its conjugate exactly when the Yukawa pattern $u=q+h$, $d=q-h$, $e=l-h$, $\nu=l+h$ holds (Lean: \texttt{SectorLinks.links\_iff\_yukawa}). For the derived charges every hop has $|\Delta Y|=1$ (\texttt{standard\_links}). The invariant-Yukawa premise of the charge derivation is therefore the gauge-covariance condition of the domain-wall hop, and $Y_H=1$ is the hop's charge; \texttt{forced\_by\_links} restates the charge selection in this form. With the two chiralities on different sectors, the flux conditions are exactly the three linear gauge-anomaly conditions (gravitational, $SU(3)^2U(1)$, $SU(2)^2U(1)$; \texttt{ChiralAnomalies.flux\_is\_linear\_anomalies}); colour is vector-like and the four doublets per generation avoid the global $SU(2)$ anomaly. The forced charges also cancel the cubic $U(1)^3$ anomaly, which was not assumed (\texttt{forced\_anomaly\_free}). The neutral sterile channel need not be assumed either. The link is the scalar, and the scalar is the regulator of the shared lepton-doublet reference (\cref{app:gauge_admittance}); like that reference it is a colourless weak doublet, and carrying the reference's hypercharge up to conjugation fixes $|Y_L|=h=1$. Equivalently, the link's components are the reference's displacement directions, with charges $0$ and $-1$; a doublet of hypercharge $h$ has charges $(h\pm1)/2$, which are $\{0,\mp1\}$ exactly when $|h|=1$ (\texttt{ReferenceMasses.link\_charge}). With $Y_L=-1$ the charges are the derived ones and $Y_\nu=0$ (\texttt{ReferenceLink.neutral\_from\_reference}); the other sign only exchanges which singlet is neutral (\texttt{conjugate\_or\_same}), and the alternative charges that \texttt{neutrality\_needed} exhibits have $Y_L=0$, a link that does not carry the reference's charge. The neutral channel is then the unique total singlet: every colour, weak and hypercharge generator fixes $\nu_R$ and moves every other carrier (\texttt{sterile\_singlet}, \texttt{sterile\_unique}). That the link carries the reference's charge is itself required for the regulation to persist: the restoring term pairs the regulator (charge $a$) with the displacement it cancels (charge $b$), changes by $e^{i\theta(a+b)}$ under a phase action, and is invariant for all $\theta$ exactly when $a=-b$ (\texttt{PersistenceClosures.invariant\_iff}). A non-invariant term would be moved by the gauge action. With $Y_L=-1$ this gives $h=1$ and the derived charges (\texttt{link\_from\_covariance}).

\paragraph{The representation from the carriers.} The channel content need not be supplied separately. The eight minimal vectors of the D$_4$ vector class are $\pm2e_i$, four axes with two signs, the $(4,2)$ of $SU(4)\times SU(2)$ with lepton number as a fourth colour (\texttt{CarrierRepresentation.one\_axis}). Singling out one axis as the lepton axis leaves six carriers on the other three and two on it, and reading the sign as the doublet index gives exactly the multiplicities $6,3,3,2,1,1$ (\texttt{multiplicities\_from\_carriers}). A traceless lepton-number charge equal on the colour axes, normalized by the forced lepton charge, is $1/3$ per colour axis, and $Y=B-L$ on the left and $B-L+2T_3$ on the right reproduce the derived charges (\texttt{charges\_from\_carriers}). A generation's two sectors must use distinct glue classes, since equal classes would identify a left and a right carrier as the same substrate vector (\texttt{Generations.same\_class\_loses}); with no sector label, generations are unordered pairs of the three nontrivial classes, of which there are exactly three (\texttt{three\_pairs}). The lepton axis is the reference's. The reserved cell holds one carrier state, and every vector-class carrier lies on exactly one axis, so the reference singles out exactly one. A Hermitian $SU(4)$ generator fixes that axis exactly when it is the extension by zero of a Hermitian generator on the other three, with the same trace (\texttt{ReferenceAxis.axis\_stabilizer}), so the unbroken part of $SU(4)$ is $SU(3)$ on the remaining three axes: colour is the stabilizer of the reference's axis, and there are $4-1=3$ colours. $B-L$ itself is moved by the reference (\texttt{bl\_broken}), as hypercharge is. Premises: triality carries each sector to the vector class, and the sign is the doublet index.
